% Source: https://artofproblemsolving.com/community/c3h2696599_nice_geometry
% Stable identity: aops-topic-2696599-q01
Let $\angle BAC=2\theta$, so $\angle BCA=90^\circ-2\theta$. Let the bisector of $\angle A$ meet $BC$ at $D$; in right triangle $ABD$, $\angle ADB=90^\circ-\theta$. Since $CP=CB$, triangle $BCP$ is isosceles with $\angle CBP=\angle CPB$. There are two positions for $P$ (called $P_1$ and $P_2$ in the sketch).

\emph{Case 1: $P$ and $B$ are on opposite sides of $AC$.} Going around $C$, the rays $CB$, $CA$, $CP$ come in this order, so $\angle BCP=(90^\circ-2\theta)+90^\circ=180^\circ-2\theta<180^\circ$, and $P$ lies on the same side of $BC$ as $A$. Hence $\angle CBP=\dfrac{180^\circ-(180^\circ-2\theta)}{2}=\theta$. The rays $BP$ and $DA$ leave line $BC$ into the same half-plane with $\angle PBD=\theta$ and $\angle ADB=90^\circ-\theta$, so they meet at a point $Q$ with $\angle BQD=180^\circ-\theta-(90^\circ-\theta)=90^\circ$. Thus $BP\perp AD$.

\emph{Case 2: $P$ and $B$ are on the same side of $AC$.} Now $\angle ACB<90^\circ=\angle ACP$, so ray $CB$ lies inside $\angle ACP$, $\angle BCP=90^\circ-(90^\circ-2\theta)=2\theta$, and $P$, $A$ are on opposite sides of line $BC$. Hence $\angle CBP=\dfrac{180^\circ-2\theta}{2}=90^\circ-\theta=\angle ADB$. These are alternate angles for the lines $BP$ and $DA$ cut by the transversal $BD$, so $BP\parallel AD$.

In both cases line $BP$ is perpendicular or parallel to the bisector of $\angle A$. $\blacksquare$

\begin{asy}
usepackage("amsmath");
size(5cm);
defaultpen(fontsize(10pt));
pair B=(0,0),A=(4,0),C=(0,3),D=(0,4/3);
pair P1=C+(3/5)*(3,4),P2=C-(3/5)*(3,4);
pair u=(A-C)/5,v=(0.6,0.8);
draw(A--B--C--cycle,linewidth(0.8));draw(C--P1);draw(C--P2);
draw(A--(A+1.45*(D-A)),blue);draw(B--P1,dashed);draw(B--P2,dashed);
draw((0.22,0)--(0.22,0.22)--(0,0.22));
draw(C+0.22*u--C+0.22*(u+v)--C+0.22*v);draw(C+0.22*u--C+0.22*(u-v)--C-0.22*v);
dot(D);
label("$A$",A,SE);label("$B$",B,SW);label("$C$",C,W);label("$D$",D,SW);
label("$P_1$",P1,N);label("$P_2$",P2,W);
\end{asy}
